\chapter{ComponentStory v0}\label{record-039}

\noindent{\footnotesize\color{BookMuted}\path{docs/strategy/component-story-v0.md}\par}\smallskip

\begin{quote}
Historical note: this is an early DOGFOOD substrate design note. Use \href{https://github.com/flyingrobots/bijou/blob/cbd70612140e7a450d764b3c7be1c6af878b99bc/docs/DOGFOOD.md}{\texttt{d\allowbreak{}o\allowbreak{}c\allowbreak{}s\allowbreak{}/\allowbreak{}D\allowbreak{}O\allowbreak{}G\allowbreak{}F\allowbreak{}O\allowbreak{}O\allowbreak{}D\allowbreak{}.\allowbreak{}m\allowbreak{}d}} for the current docs-surface signpost and METHOD backlog/design docs for active queue work.
\end{quote}

\emph{Design note for the first DOGFOOD substrate}

\section{Why this exists}\label{record-039-why-this-exists}

Bijou currently has several near-miss representations of ``a component in context'':

\begin{itemize}
\tightlist
\item
  design-system family docs
\item
  package/example README prose
\item
  showcase \texttt{C\allowbreak{}o\allowbreak{}m\allowbreak{}p\allowbreak{}o\allowbreak{}n\allowbreak{}e\allowbreak{}n\allowbreak{}t\allowbreak{}E\allowbreak{}n\allowbreak{}t\allowbreak{}r\allowbreak{}y}
\item
  smoke and canary examples
\end{itemize}

Those artifacts overlap, but they do not share one canonical contract. \texttt{C\allowbreak{}o\allowbreak{}m\allowbreak{}p\allowbreak{}o\allowbreak{}n\allowbreak{}e\allowbreak{}n\allowbreak{}t\allowbreak{}S\allowbreak{}t\allowbreak{}o\allowbreak{}r\allowbreak{}y\allowbreak{} \allowbreak{}v\allowbreak{}0} is the first attempt to fix that.

The goal is not to design the forever-schema. The goal is to create the smallest structured record that can prove this loop:

\begin{enumerate}
\def\labelenumi{\arabic{enumi}.}
\tightlist
\item
  a component family can be documented from structured fields
\item
  the same record can render a live demo
\item
  the same record can show graceful lowering across profiles
\item
  the same record can later grow selectors and scenarios without being thrown away
\end{enumerate}

\section{\texorpdfstring{The problem with \texttt{C\allowbreak{}o\allowbreak{}m\allowbreak{}p\allowbreak{}o\allowbreak{}n\allowbreak{}e\allowbreak{}n\allowbreak{}t\allowbreak{}E\allowbreak{}n\allowbreak{}t\allowbreak{}r\allowbreak{}y}}{The problem with }}\label{record-039-the-problem-with-componententry}

Today's showcase contract in \href{https://github.com/flyingrobots/bijou/blob/cbd70612140e7a450d764b3c7be1c6af878b99bc/examples/showcase/types.ts}{examples/showcase/types.ts} is useful, but too compressed:

\begin{itemize}
\tightlist
\item
  \texttt{d\allowbreak{}e\allowbreak{}s\allowbreak{}c\allowbreak{}r\allowbreak{}i\allowbreak{}p\allowbreak{}t\allowbreak{}i\allowbreak{}o\allowbreak{}n} is one markdown blob
\item
  \texttt{r\allowbreak{}e\allowbreak{}n\allowbreak{}d\allowbreak{}e\allowbreak{}r\allowbreak{}(\allowbreak{}w\allowbreak{}i\allowbreak{}d\allowbreak{}t\allowbreak{}h\allowbreak{},\allowbreak{} \allowbreak{}c\allowbreak{}t\allowbreak{}x\allowbreak{})} is one opaque demo function
\item
  graceful lowering is mostly described in prose, not structured
\item
  ``when to use / when not to use'' is not first-class data
\item
  selectors, scenarios, source snippets, and variant controls have nowhere to live
\end{itemize}

That means it is good enough for a showcase, but not strong enough to become the substrate for docs, tests, and tooling.

\section{Design goals}\label{record-039-design-goals}

\texttt{C\allowbreak{}o\allowbreak{}m\allowbreak{}p\allowbreak{}o\allowbreak{}n\allowbreak{}e\allowbreak{}n\allowbreak{}t\allowbreak{}S\allowbreak{}t\allowbreak{}o\allowbreak{}r\allowbreak{}y\allowbreak{} \allowbreak{}v\allowbreak{}0} should:

\begin{enumerate}
\def\labelenumi{\arabic{enumi}.}
\item
  \textbf{Be structured enough for docs rendering} The docs surface should not have to parse one giant markdown blob to find summary, use, avoid, and lowering guidance.
\item
  \textbf{Stay small enough to author by hand} If the schema is too heavy, maintainers will bypass it and drift returns immediately.
\item
  \textbf{Support profile-aware demos directly} A story should be able to render itself in \texttt{r\allowbreak{}i\allowbreak{}c\allowbreak{}h}, \texttt{s\allowbreak{}t\allowbreak{}a\allowbreak{}t\allowbreak{}i\allowbreak{}c}, \texttt{p\allowbreak{}i\allowbreak{}p\allowbreak{}e}, and \texttt{a\allowbreak{}c\allowbreak{}c\allowbreak{}e\allowbreak{}s\allowbreak{}s\allowbreak{}i\allowbreak{}b\allowbreak{}l\allowbreak{}e}.
\item
  \textbf{Model variants explicitly} Component families often have more than one canonical configuration worth teaching.
\item
  \textbf{Leave room for selectors and scenarios later} v0 should not depend on a selector model or scenario harness, but it should not block them either.
\end{enumerate}

\section{Non-goals for v0}\label{record-039-non-goals-for-v0}

\texttt{C\allowbreak{}o\allowbreak{}m\allowbreak{}p\allowbreak{}o\allowbreak{}n\allowbreak{}e\allowbreak{}n\allowbreak{}t\allowbreak{}S\allowbreak{}t\allowbreak{}o\allowbreak{}r\allowbreak{}y\allowbreak{} \allowbreak{}v\allowbreak{}0} is not trying to solve:

\begin{itemize}
\tightlist
\item
  every possible interactive scenario
\item
  full \texttt{P\allowbreak{}a\allowbreak{}t\allowbreak{}t\allowbreak{}e\allowbreak{}r\allowbreak{}n\allowbreak{}S\allowbreak{}t\allowbreak{}o\allowbreak{}r\allowbreak{}y} design
\item
  visual editor compatibility
\item
  complete replay/assertion semantics
\item
  an exhaustive control/knob system
\item
  every component in the repo at once
\end{itemize}

\section{Proposed shape}\label{record-039-proposed-shape}

\begin{CodeBlock}
type StoryPackage = 'bijou' | 'bijou-tui' | 'bijou-tui-app';
type StoryMode = 'interactive' | 'static' | 'pipe' | 'accessible';

type StoryPreview = string | Surface;

interface StoryLowering {
  readonly interactive: string;
  readonly static: string;
  readonly pipe: string;
  readonly accessible: string;
}

interface StoryDocs {
  readonly summary: string;
  readonly useWhen: readonly string[];
  readonly avoidWhen: readonly string[];
  readonly relatedFamilies: readonly string[];
  readonly gracefulLowering: StoryLowering;
}

interface StoryVariant<State = void> {
  readonly id: string;
  readonly label: string;
  readonly description?: string;
  readonly initialState?: State;
  readonly render: (input: {
    readonly width: number;
    readonly ctx: BijouContext;
    readonly state: State;
  }) => StoryPreview;
}

interface StoryProfilePreset {
  readonly id: string;
  readonly label: string;
  readonly mode: StoryMode;
  readonly width: number;
}

interface StorySource {
  readonly examplePath?: string;
  readonly snippetLabel?: string;
}

interface ComponentStory<State = void> {
  readonly kind: 'component';
  readonly id: string;
  readonly family: string;
  readonly title: string;
  readonly package: StoryPackage;
  readonly docs: StoryDocs;
  readonly profilePresets: readonly StoryProfilePreset[];
  readonly variants: readonly StoryVariant<State>[];
  readonly source?: StorySource;
  readonly tags?: readonly string[];
}
\end{CodeBlock}

\section{Why this shape}\label{record-039-why-this-shape}

\subsection{\texorpdfstring{\texttt{d\allowbreak{}o\allowbreak{}c\allowbreak{}s}}{}}\label{record-039-docs}

This is the biggest difference from \texttt{C\allowbreak{}o\allowbreak{}m\allowbreak{}p\allowbreak{}o\allowbreak{}n\allowbreak{}e\allowbreak{}n\allowbreak{}t\allowbreak{}E\allowbreak{}n\allowbreak{}t\allowbreak{}r\allowbreak{}y}.

Instead of one markdown blob, the story has:

\begin{itemize}
\tightlist
\item
  \texttt{s\allowbreak{}u\allowbreak{}m\allowbreak{}m\allowbreak{}a\allowbreak{}r\allowbreak{}y}
\item
  \texttt{u\allowbreak{}s\allowbreak{}e\allowbreak{}W\allowbreak{}h\allowbreak{}e\allowbreak{}n}
\item
  \texttt{a\allowbreak{}v\allowbreak{}o\allowbreak{}i\allowbreak{}d\allowbreak{}W\allowbreak{}h\allowbreak{}e\allowbreak{}n}
\item
  \texttt{r\allowbreak{}e\allowbreak{}l\allowbreak{}a\allowbreak{}t\allowbreak{}e\allowbreak{}d\allowbreak{}F\allowbreak{}a\allowbreak{}m\allowbreak{}i\allowbreak{}l\allowbreak{}i\allowbreak{}e\allowbreak{}s}
\item
  \texttt{g\allowbreak{}r\allowbreak{}a\allowbreak{}c\allowbreak{}e\allowbreak{}f\allowbreak{}u\allowbreak{}l\allowbreak{}L\allowbreak{}o\allowbreak{}w\allowbreak{}e\allowbreak{}r\allowbreak{}i\allowbreak{}n\allowbreak{}g}
\end{itemize}

That makes the docs surface renderable without brittle parsing and keeps the design-system doctrine aligned with the story itself.

\subsection{\texorpdfstring{\texttt{p\allowbreak{}r\allowbreak{}o\allowbreak{}f\allowbreak{}i\allowbreak{}l\allowbreak{}e\allowbreak{}P\allowbreak{}r\allowbreak{}e\allowbreak{}s\allowbreak{}e\allowbreak{}t\allowbreak{}s}}{}}\label{record-039-profilepresets}

This makes profile comparison explicit and small.

v0 should not model every possible runtime knob. It only needs enough to render canonical comparisons:

\begin{itemize}
\tightlist
\item
  one rich preset
\item
  one static preset
\item
  one pipe preset
\item
  one accessible preset
\end{itemize}

Width is included because graceful lowering is often width-sensitive in terminal UIs.

\subsection{\texorpdfstring{\texttt{v\allowbreak{}a\allowbreak{}r\allowbreak{}i\allowbreak{}a\allowbreak{}n\allowbreak{}t\allowbreak{}s}}{}}\label{record-039-variants}

Variants are first-class because many families are not taught well with one demo:

\begin{itemize}
\tightlist
\item
  \texttt{a\allowbreak{}l\allowbreak{}e\allowbreak{}r\allowbreak{}t\allowbreak{}(\allowbreak{})} needs more than one tone
\item
  \texttt{m\allowbreak{}o\allowbreak{}d\allowbreak{}a\allowbreak{}l\allowbreak{}(\allowbreak{})} may need a simple and structured-body variant
\item
  \texttt{v\allowbreak{}i\allowbreak{}e\allowbreak{}w\allowbreak{}p\allowbreak{}o\allowbreak{}r\allowbreak{}t\allowbreak{}S\allowbreak{}u\allowbreak{}r\allowbreak{}f\allowbreak{}a\allowbreak{}c\allowbreak{}e\allowbreak{}(\allowbreak{})} may need one text case and one structured case
\end{itemize}

v0 keeps this minimal:

\begin{itemize}
\tightlist
\item
  label
\item
  optional description
\item
  optional initial state
\item
  render function
\end{itemize}

\subsection{\texorpdfstring{\texttt{s\allowbreak{}o\allowbreak{}u\allowbreak{}r\allowbreak{}c\allowbreak{}e}}{}}\label{record-039-source}

This stays intentionally light in v0.

The docs app only needs a way to say:

\begin{itemize}
\tightlist
\item
  where the canonical example lives
\item
  or what snippet label should be shown
\end{itemize}

Anything more ambitious can come later with story tooling.

\section{\texorpdfstring{Migration from \texttt{C\allowbreak{}o\allowbreak{}m\allowbreak{}p\allowbreak{}o\allowbreak{}n\allowbreak{}e\allowbreak{}n\allowbreak{}t\allowbreak{}E\allowbreak{}n\allowbreak{}t\allowbreak{}r\allowbreak{}y}}{Migration from }}\label{record-039-migration-from-componententry}

The intended mapping is:

{\def\LTcaptype{none} % do not increment counter
\begin{longtable}[]{@{}
  >{\raggedright\arraybackslash}p{(\linewidth - 2\tabcolsep) * \real{0.5000}}
  >{\raggedright\arraybackslash}p{(\linewidth - 2\tabcolsep) * \real{0.5000}}@{}}
\toprule\noalign{}
\texttt{C\allowbreak{}o\allowbreak{}m\allowbreak{}p\allowbreak{}o\allowbreak{}n\allowbreak{}e\allowbreak{}n\allowbreak{}t\allowbreak{}E\allowbreak{}n\allowbreak{}t\allowbreak{}r\allowbreak{}y}
 & \texttt{C\allowbreak{}o\allowbreak{}m\allowbreak{}p\allowbreak{}o\allowbreak{}n\allowbreak{}e\allowbreak{}n\allowbreak{}t\allowbreak{}S\allowbreak{}t\allowbreak{}o\allowbreak{}r\allowbreak{}y\allowbreak{} \allowbreak{}v\allowbreak{}0}
 \\
\midrule\noalign{}
\endhead
\bottomrule\noalign{}
\endlastfoot
\texttt{i\allowbreak{}d} & \texttt{i\allowbreak{}d} \\
\texttt{n\allowbreak{}a\allowbreak{}m\allowbreak{}e} & \texttt{t\allowbreak{}i\allowbreak{}t\allowbreak{}l\allowbreak{}e} \\
\texttt{s\allowbreak{}u\allowbreak{}b\allowbreak{}t\allowbreak{}i\allowbreak{}t\allowbreak{}l\allowbreak{}e} & \texttt{d\allowbreak{}o\allowbreak{}c\allowbreak{}s\allowbreak{}.\allowbreak{}s\allowbreak{}u\allowbreak{}m\allowbreak{}m\allowbreak{}a\allowbreak{}r\allowbreak{}y} \\
\texttt{d\allowbreak{}e\allowbreak{}s\allowbreak{}c\allowbreak{}r\allowbreak{}i\allowbreak{}p\allowbreak{}t\allowbreak{}i\allowbreak{}o\allowbreak{}n} & {s\allowbreak{}p\allowbreak{}l\allowbreak{}i\allowbreak{}t} {i\allowbreak{}n\allowbreak{}t\allowbreak{}o} \texttt{d\allowbreak{}o\allowbreak{}c\allowbreak{}s\allowbreak{}.\allowbreak{}*} {f\allowbreak{}i\allowbreak{}e\allowbreak{}l\allowbreak{}d\allowbreak{}s} \\
\texttt{p\allowbreak{}k\allowbreak{}g} & \texttt{p\allowbreak{}a\allowbreak{}c\allowbreak{}k\allowbreak{}a\allowbreak{}g\allowbreak{}e} \\
\texttt{r\allowbreak{}e\allowbreak{}n\allowbreak{}d\allowbreak{}e\allowbreak{}r\allowbreak{}(\allowbreak{}w\allowbreak{}i\allowbreak{}d\allowbreak{}t\allowbreak{}h\allowbreak{},\allowbreak{} \allowbreak{}c\allowbreak{}t\allowbreak{}x\allowbreak{})} & \texttt{v\allowbreak{}a\allowbreak{}r\allowbreak{}i\allowbreak{}a\allowbreak{}n\allowbreak{}t\allowbreak{}s\allowbreak{}[\allowbreak{}0\allowbreak{}]\allowbreak{}.\allowbreak{}r\allowbreak{}e\allowbreak{}n\allowbreak{}d\allowbreak{}e\allowbreak{}r\allowbreak{}(\allowbreak{}.\allowbreak{}.\allowbreak{}.\allowbreak{})} \\
\texttt{t\allowbreak{}i\allowbreak{}e\allowbreak{}r} & {d\allowbreak{}r\allowbreak{}o\allowbreak{}p\allowbreak{}p\allowbreak{}e\allowbreak{}d} {f\allowbreak{}r\allowbreak{}o\allowbreak{}m} {v\allowbreak{}0} {s\allowbreak{}t\allowbreak{}o\allowbreak{}r\allowbreak{}y} {c\allowbreak{}o\allowbreak{}r\allowbreak{}e\allowbreak{};} {t\allowbreak{}h\allowbreak{}i\allowbreak{}s} {i\allowbreak{}s} {s\allowbreak{}h\allowbreak{}o\allowbreak{}w\allowbreak{}c\allowbreak{}a\allowbreak{}s\allowbreak{}e\allowbreak{}/\allowbreak{}w\allowbreak{}o\allowbreak{}r\allowbreak{}k\allowbreak{}b\allowbreak{}e\allowbreak{}n\allowbreak{}c\allowbreak{}h} {m\allowbreak{}e\allowbreak{}t\allowbreak{}a\allowbreak{}d\allowbreak{}a\allowbreak{}t\allowbreak{}a\allowbreak{},} {n\allowbreak{}o\allowbreak{}t} {s\allowbreak{}t\allowbreak{}o\allowbreak{}r\allowbreak{}y} {t\allowbreak{}r\allowbreak{}u\allowbreak{}t\allowbreak{}h} \\
\end{longtable}
}

The important design choice here is dropping \texttt{t\allowbreak{}i\allowbreak{}e\allowbreak{}r} from the canonical story model.

Why:

\begin{itemize}
\tightlist
\item
  \texttt{t\allowbreak{}i\allowbreak{}e\allowbreak{}r} is about embedding/showcase ergonomics
\item
  it is not a core truth about the component family itself
\item
  if a future workbench needs it, it can layer that on top instead of contaminating the story substrate
\end{itemize}

\section{Recommended first stories}\label{record-039-recommended-first-stories}

The first v0 stories should cover three different teaching problems:

\subsection{\texorpdfstring{1. \texttt{a\allowbreak{}l\allowbreak{}e\allowbreak{}r\allowbreak{}t\allowbreak{}(\allowbreak{})} / \texttt{n\allowbreak{}o\allowbreak{}t\allowbreak{}e\allowbreak{}(\allowbreak{})}}{1.  / }}\label{record-039-alert-note}

Why:

\begin{itemize}
\tightlist
\item
  low interaction complexity
\item
  clear graceful lowering
\item
  easy to render side-by-side
\end{itemize}

What they validate:

\begin{itemize}
\tightlist
\item
  structured docs fields
\item
  simple profile comparison
\item
  variant support
\end{itemize}

\subsection{\texorpdfstring{2. \texttt{m\allowbreak{}o\allowbreak{}d\allowbreak{}a\allowbreak{}l\allowbreak{}(\allowbreak{})} / \texttt{d\allowbreak{}r\allowbreak{}a\allowbreak{}w\allowbreak{}e\allowbreak{}r\allowbreak{}(\allowbreak{})}}{2.  / }}\label{record-039-modal-drawer}

Why:

\begin{itemize}
\tightlist
\item
  structured overlay content
\item
  clearer ``use vs avoid'' doctrine
\item
  richer TUI-only behavior without requiring a full app shell
\end{itemize}

What they validate:

\begin{itemize}
\tightlist
\item
  surface-native story previews
\item
  overlay-family doctrine in structured form
\item
  more than one meaningful variant
\end{itemize}

\subsection{\texorpdfstring{3. \texttt{v\allowbreak{}i\allowbreak{}e\allowbreak{}w\allowbreak{}p\allowbreak{}o\allowbreak{}r\allowbreak{}t\allowbreak{}S\allowbreak{}u\allowbreak{}r\allowbreak{}f\allowbreak{}a\allowbreak{}c\allowbreak{}e\allowbreak{}(\allowbreak{})} or \texttt{n\allowbreak{}a\allowbreak{}v\allowbreak{}i\allowbreak{}g\allowbreak{}a\allowbreak{}b\allowbreak{}l\allowbreak{}e\allowbreak{}T\allowbreak{}a\allowbreak{}b\allowbreak{}l\allowbreak{}e\allowbreak{}S\allowbreak{}u\allowbreak{}r\allowbreak{}f\allowbreak{}a\allowbreak{}c\allowbreak{}e\allowbreak{}(\allowbreak{})}}{3.  or }}\label{record-039-viewportsurface-or-navigabletablesurface}

Why:

\begin{itemize}
\tightlist
\item
  tests a component whose value is not just appearance but behavior and composition
\end{itemize}

What they validate:

\begin{itemize}
\tightlist
\item
  structured TUI docs can still describe behavior honestly
\item
  the docs surface can present width-sensitive behavior
\item
  later selector/scenario integration has a real target
\end{itemize}

\section{Open questions}\label{record-039-open-questions}

These are intentionally deferred, but should stay visible:

\begin{enumerate}
\def\labelenumi{\arabic{enumi}.}
\tightlist
\item
  Should \texttt{C\allowbreak{}o\allowbreak{}m\allowbreak{}p\allowbreak{}o\allowbreak{}n\allowbreak{}e\allowbreak{}n\allowbreak{}t\allowbreak{}S\allowbreak{}t\allowbreak{}o\allowbreak{}r\allowbreak{}y} carry source snippets directly, or only references to example files?
\item
  Should \texttt{f\allowbreak{}a\allowbreak{}m\allowbreak{}i\allowbreak{}l\allowbreak{}y} be a free string in v0, or a stricter id shared with the design-system docs?
\item
  Should profile presets live on the story, or should the docs app own the canonical presets globally?
\item
  Should \texttt{P\allowbreak{}a\allowbreak{}t\allowbreak{}t\allowbreak{}e\allowbreak{}r\allowbreak{}n\allowbreak{}S\allowbreak{}t\allowbreak{}o\allowbreak{}r\allowbreak{}y} share a common base with \texttt{C\allowbreak{}o\allowbreak{}m\allowbreak{}p\allowbreak{}o\allowbreak{}n\allowbreak{}e\allowbreak{}n\allowbreak{}t\allowbreak{}S\allowbreak{}t\allowbreak{}o\allowbreak{}r\allowbreak{}y}, or be defined independently later?
\end{enumerate}

\section{Recommended implementation sequence}\label{record-039-recommended-implementation-sequence}

\begin{enumerate}
\def\labelenumi{\arabic{enumi}.}
\tightlist
\item
  Define the v0 type in code without promising permanence.
\item
  Author three canonical stories by hand.
\item
  Build the first docs surface around those stories.
\item
  Only after that, revisit selectors and scenarios.
\end{enumerate}

That order matters because the immediate question is not ``what is the perfect schema?'' It is ``can one structured story record actually replace duplicated docs/demo truth?''

